\documentclass[10pt,a4paper]{amsart}
\usepackage{amsmath,amssymb,amsthm,graphicx}
\usepackage{url}
\usepackage{hyperref}
\usepackage{verbatim}
\newenvironment{reference}{\comment}{\endcomment}
%\newenvironment{reference}{}{}
\newenvironment{slogan}{\comment}{\endcomment}
\newenvironment{history}{\comment}{\endcomment}

% For commutative diagrams we use Xy-pic
\usepackage[all]{xy}

% We use 2cell for 2-commutative diagrams.
\xyoption{2cell}
% Theorem environments.
%
\theoremstyle{plain}
\newtheorem{theorem}[subsection]{Theorem}
\newtheorem{proposition}[subsection]{Proposition}
\newtheorem{lemma}[subsection]{Lemma}

\theoremstyle{definition}
\newtheorem{definition}[subsection]{Definition}
\newtheorem{example}[subsection]{Example}
\newtheorem{exercise}[subsection]{Exercise}
\newtheorem{situation}[subsection]{Situation}

\theoremstyle{remark}
\newtheorem{remark}[subsection]{Remark}
\newtheorem{remarks}[subsection]{Remarks}

\numberwithin{equation}{subsection}

% Macros
%
\def\lim{\mathop{\mathrm{lim}}\nolimits}
\def\colim{\mathop{\mathrm{colim}}\nolimits}
\def\Spec{\mathop{\mathrm{Spec}}}
\def\Hom{\mathop{\mathrm{Hom}}\nolimits}
\def\Ext{\mathop{\mathrm{Ext}}\nolimits}
\def\SheafHom{\mathop{\mathcal{H}\!\mathit{om}}\nolimits}
\def\SheafExt{\mathop{\mathcal{E}\!\mathit{xt}}\nolimits}
\def\Sch{\mathit{Sch}}
\def\Mor{\mathop{\mathrm{Mor}}\nolimits}
\def\Ob{\mathop{\mathrm{Ob}}\nolimits}
\def\Sh{\mathop{\mathit{Sh}}\nolimits}
\def\NL{\mathop{N\!L}\nolimits}
\def\CH{\mathop{\mathrm{CH}}\nolimits}
\def\proetale{{pro\text{-}\acute{e}tale}}
\def\etale{{\acute{e}tale}}
\def\QCoh{\mathit{QCoh}}
\def\Ker{\mathop{\mathrm{Ker}}}
\def\Im{\mathop{\mathrm{Im}}}
\def\Coker{\mathop{\mathrm{Coker}}}
\def\Coim{\mathop{\mathrm{Coim}}}

% Boxtimes
%
\DeclareMathSymbol{\boxtimes}{\mathbin}{AMSa}{"02}

%
% Macros for moduli stacks/spaces
%
\def\QCohstack{\mathcal{QC}\!\mathit{oh}}
\def\Cohstack{\mathcal{C}\!\mathit{oh}}
\def\Spacesstack{\mathcal{S}\!\mathit{paces}}
\def\Quotfunctor{\mathrm{Quot}}
\def\Hilbfunctor{\mathrm{Hilb}}
\def\Curvesstack{\mathcal{C}\!\mathit{urves}}
\def\Polarizedstack{\mathcal{P}\!\mathit{olarized}}
\def\Complexesstack{\mathcal{C}\!\mathit{omplexes}}
% \Pic is the operator that assigns to X its picard group, usage \Pic(X)
% \Picardstack_{X/B} denotes the Picard stack of X over B
% \Picardfunctor_{X/B} denotes the Picard functor of X over B
\def\Pic{\mathop{\mathrm{Pic}}\nolimits}
\def\Picardstack{\mathcal{P}\!\mathit{ic}}
\def\Picardfunctor{\mathrm{Pic}}
\def\Deformationcategory{\mathcal{D}\!\mathit{ef}}
\usepackage{cleveref}
\crefname{originaldefinition}{original definition}{original definitions}
\Crefname{originaldefinition}{Original definition}{Original definitions}
\makeatletter
\newlabel{algebra-definition-almost-integral}{{00GW}{source}{}{sourcepointer-1}{}}
\newlabel{algebra-definition-almost-integral@cref}{{[originaldefinition][0][]00GW}{[1][0][]source}{}{}{}}
\makeatother
\crefname{originalsection}{original section}{original sections}
\Crefname{originalsection}{Original section}{Original sections}
\makeatletter
\newlabel{algebra-example-valuation-ring-coherent}{{0EWV}{source}{}{sourcepointer-2}{}}
\newlabel{algebra-example-valuation-ring-coherent@cref}{{[originalsection][0][]0EWV}{[1][0][]source}{}{}{}}
\makeatother
\crefname{originallemma}{original lemma}{original lemmas}
\Crefname{originallemma}{Original lemma}{Original lemmas}
\makeatletter
\newlabel{algebra-lemma-ff-rings}{{00HQ}{source}{}{sourcepointer-3}{}}
\newlabel{algebra-lemma-ff-rings@cref}{{[originallemma][0][]00HQ}{[1][0][]source}{}{}{}}
\makeatother
\crefname{originallemma}{original lemma}{original lemmas}
\Crefname{originallemma}{Original lemma}{Original lemmas}
\makeatletter
\newlabel{algebra-lemma-valuation-ring-normal}{{00IC}{source}{}{sourcepointer-4}{}}
\newlabel{algebra-lemma-valuation-ring-normal@cref}{{[originallemma][0][]00IC}{[1][0][]source}{}{}{}}
\makeatother
\crefname{originalproposition}{original proposition}{original propositions}
\Crefname{originalproposition}{Original proposition}{Original propositions}
\makeatletter
\newlabel{algebra-proposition-characterize-coherent}{{05CZ}{source}{}{sourcepointer-5}{}}
\newlabel{algebra-proposition-characterize-coherent@cref}{{[originalproposition][0][]05CZ}{[1][0][]source}{}{}{}}
\makeatother
\title{For a valuation ring of dimension > 1, R[[x]] is flat over R but not flat over R[x]}
\author{The Stacks Project Authors}
\date{stacks-project:examples.tex:local-snapshot-20261012; GNU FDL 1.2 or later}
\begin{document}
\maketitle
\noindent\textit{Editorial note.} All delivered mathematical text is the author's own, copied from the pinned Stacks Project chapter file examples.tex; no statement or proof step was rewritten, reordered or removed.\par\medskip
\noindent\textit{Editorial note.} The counted proof is the author's complete argument: The opening sentence 'The arguments above show that this is true if we can show that R is not completely normal' relies on the context span; the reduction and the almost-integral construction of 1/y are both inside the counted span. Flatness over R itself comes from the context paragraph (valuation rings coherent), completing the flat-over-R / not-flat-over-R[x] pair.\par\medskip
\noindent\textit{Editorial note.} The label lemma-completion-polynomial-ring-not-flat-bis names the counted statement, and the proof environment immediately following it in the source is delivered verbatim as its proof.\par\medskip
\noindent\textit{Editorial note.} The counted statement is carried by the author's own lemma environment in the chapter source.\par\medskip
\noindent\textit{Allowed standard prerequisites (editorial).}\begin{itemize}
\item Context span 1042-1083: Badam-Baplan construction - a valuation ring R is coherent (Algebra, example-valuation-ring-coherent) so R[[x]] is flat over R (Algebra, proposition-characterize-coherent); and the same-section lemma proving: if R[[x]] is flat over R[x] for a domain R with fraction field K, then R is normal iff completely normal (almost-integral elements detected via (ax-b)f = -rb and the faithfully flat localization at S = \{h : h(0)=1\})
\item Valuation rings are normal (Algebra, lemma-valuation-ring-normal); almost integral elements: ralpha\textasciicircum{}n in R for all n
\end{itemize}\medskip
\ifdefined\hypertarget\hypertarget{sourcepointer-1}{}\fi
\noindent\textit{Source reference.} Original definition 00GW: sources/stacks-chow/tags.txt, source lines 622--622. Prerequisite pointer invoked at 1117; it names a standard fact of a sibling chapter that the counted argument uses, and every step of the counted proof built on it is retained verbatim in the delivered body.\par\medskip
\ifdefined\hypertarget\hypertarget{sourcepointer-2}{}\fi
\noindent\textit{Source reference.} Original section 0EWV: sources/stacks-chow/tags.txt, source lines 18635--18635. Prerequisite pointer invoked at 1107; it names a standard fact of a sibling chapter that the counted argument uses, and every step of the counted proof built on it is retained verbatim in the delivered body.\par\medskip
\ifdefined\hypertarget\hypertarget{sourcepointer-3}{}\fi
\noindent\textit{Source reference.} Original lemma 00HQ: sources/stacks-chow/tags.txt, source lines 652--652. Prerequisite pointer invoked at 1129; it names a standard fact of a sibling chapter that the counted argument uses, and every step of the counted proof built on it is retained verbatim in the delivered body.\par\medskip
\ifdefined\hypertarget\hypertarget{sourcepointer-4}{}\fi
\noindent\textit{Source reference.} Original lemma 00IC: sources/stacks-chow/tags.txt, source lines 674--674. Prerequisite pointer invoked at 1094; it names a standard fact of a sibling chapter that the counted argument uses, and every step of the counted proof built on it is retained verbatim in the delivered body.\par\medskip
\ifdefined\hypertarget\hypertarget{sourcepointer-5}{}\fi
\noindent\textit{Source reference.} Original proposition 05CZ: sources/stacks-chow/tags.txt, source lines 6937--6937. Prerequisite pointer invoked at 1109; it names a standard fact of a sibling chapter that the counted argument uses, and every step of the counted proof built on it is retained verbatim in the delivered body.\par\medskip
\noindent
It turns out there is a ring $R$ such that $R[[x]]$ is flat
over $R$, but $R[[x]]$ is not flat over $R[x]$. See
\href{https://math.stackexchange.com/users/164860/badam-baplan}{this post}
by Badam Baplan. Namely, let
$R$ be a valuation ring. Then $R$ is coherent (Algebra, Example
\ref{algebra-example-valuation-ring-coherent}) and hence
$R[[x]]$ is flat over $R$ by 
Algebra, Proposition \ref{algebra-proposition-characterize-coherent}.
On the other hand, we have the following lemma.

\begin{lemma}
\label{lemma-almost-integral-when-powerseries-flat}
Let $R$ be a domain with fraction field $K$.
If $R[[x]]$ is flat over $R[x]$, then $R$ is normal if and only
if $R$ is completely normal
(Algebra, Definition \ref{algebra-definition-almost-integral}).
\end{lemma}

\begin{proof}
Suppose we have $\alpha \in K$ and a nonzero $r \in R$ such that
$r \alpha^n \in R$ for all $n \geq 1$. Then we consider
$f = \sum r \alpha^{n - 1} x^n$ in $R[[x]]$. Write $\alpha = a/b$
for $a, b \in R$ with $b$ nonzero. Then we see that $(a x - b)f = -rb$.
It follows that $rb$ is in the ideal $(ax - b)R[[x]]$.
Let $S = \{h \in R[x] : h(0) = 1\}$. This is a multiplicative subset
and flatness of $R[x] \to R[[x]]$ implies that $S^{-1}R[x] \to R[[x]]$
is faithfully flat (details omitted; hint: use Algebra, Lemma
\ref{algebra-lemma-ff-rings}). Hence
$$
S^{-1}R[x]/(ax - b)S^{-1}R[x] \to R[[x]]/(ax - b)R[[x]]
$$
is injective. We conclude that
$h rb = (ax - b) g$ for some $h \in S$ and $g \in R[x]$.
Writing $h = 1 + h_1 x + \ldots + h_d x^d$ shows that we obtain
$$
1 + h_1 x + \ldots + h_d x^d = (1/r)(\alpha x - 1)g
$$
This factorization in $K[x]$ gives a corresponding factorization
in $K[x^{-1}]$ which shows that $\alpha$ is the root of a monic
polynomial with coefficients in $R$ as desired.
\end{proof}

\begin{lemma}
\label{lemma-completion-polynomial-ring-not-flat-bis}
If $R$ is a valuation ring of dimension $> 1$, then $R[[x]]$
is flat over $R$ but not flat over $R[x]$.
\end{lemma}

\begin{proof}
The arguments above show that this is true if we can show that
$R$ is not completely normal (valuation rings are
normal, see Algebra, Lemma \ref{algebra-lemma-valuation-ring-normal}).
Let $\mathfrak p \subset \mathfrak m \subset R$ be a chain of primes.
Pick nonzero $x \in \mathfrak p$ and $y \in \mathfrak m \setminus \mathfrak p$.
Then $x y^{-n} \in R$ for all $n \geq 1$ (if not then $y^n/x \in R$
which is absurd because $y \not \in \mathfrak p$). Hence $1/y$ is
almost integral over $R$ but not in $R$.
\end{proof}

\end{document}
